\section{The fixed-block-count grove pairing}

Let \(1\le b\le n\), and index the rows of \(H_{n,b}\) by the
noncrossing partitions \(\pi\) with \(b\) blocks and its columns by
those \(\rho\) with \(n+1-b\) blocks. Put \(H_{n,b}(\pi,\rho)=1\)
if \(\pi\vee\rho\) is the one-block partition and \(0\) otherwise.
These square matrices are the complementary blocks of the grove
pairing of Kenyon and Wilson
(\href{https://arxiv.org/abs/math/0608422}{\emph{Boundary partitions
in trees and dimers}}, \S2.2, equation (3)).

\begin{corollary}\label{cor:grove-block}
The positive Smith factors of \(H_{n,b}\) are
\[
 \operatorname{lcm}\bigl(1,\ldots,\lfloor(h(P)+1)/2\rfloor\bigr)
\]
over the Dyck paths \(P\) of length \(2n\) with exactly \(b\) upsteps
in odd-numbered positions, counting the first step as position one.
In particular \(\rk_{\F_2}H_{n,b}=\binom{n-1}{b-1}\).
\end{corollary}

\begin{proof}
Write \(a=|\pi\vee\rho|\) and \(a'=|\pi'\vee\rho'|\).
The plane bipartite multigraph with one vertex for each block of
\(\pi\) and \(\rho\), and one edge per label joining its two blocks,
has \(a\) components and \(a'\) faces. To see the latter, draw the two
families of stars in opposite hemispheres. Their complementary regions
join the white boundary intervals exactly as in
Section~\ref{sec:two-colour}. Euler's formula gives
\[
             a-a'=|\pi|+|\rho|-n-1 .
\]
Grade \(\ZZ[X,Y]\) by \(\deg X=1,\deg Y=-1\). The entries of
\(J_n/(XY)\) have the degrees on the right, and all diagonal factors
\(b_{h(P)}/(XY)\) have degree zero. In the row relation of
Theorem~\ref{thm:two-colour}, retain from the coefficient of each
smaller row \(\tau\) only its homogeneous part of degree
\(|\pi|-|\tau|\). Coordinatewise this extracts the homogeneous
part of the degree of the original entry. The remainder is still
divisible by the same homogeneous diagonal factor. We thus obtain
\[
 L\,(J_n/(XY))=D V,\qquad D_{P,P}=b_{h(P)}(X,Y)/(XY),
\]
where \(L\) is unitriangular by area and
\(\deg V_{\pi,\rho}=|\pi|+|\rho|-n-1\) whenever the entry is nonzero.
The determinant argument gives \(\det V=1\) for these choices as well.

At \(X=Y=0\) a coordinate of \(J_n/(XY)\) survives precisely when
\(a=a'=1\). The displayed degree identity identifies its complementary
blocks with \(H_{n,b}\). Only degree-zero entries of \(L,V\) can
survive, so \(L(0,0)\) preserves block count and \(V(0,0)\) pairs
complementary counts. Each complementary square block of \(V(0,0)\)
has determinant \(\pm1\), since their product up to sign is a unit.
Each block of \(L(0,0)\) is unitriangular. Finally
\(\lvert D_{P,P}(0,0)\rvert\) is the least common multiple in the
statement, by \(B_d(0)=\Psi_d(0)\) for \(d>2\) and
Corollary~\ref{integer}. In the fattening convention of
Section~\ref{sec:two-colour}, the step at endpoint \(2i\), numbered
\(2i+1\), opens an arc exactly when \(i\) is the minimum of its
block. This identifies the required paths.
\end{proof}

More generally, put \(u=p^e\) for a prime \(p\) and \(e\ge1\).
Let \(T_u(z)\) be the \((u-1)\)-by-\((u-1)\) tridiagonal matrix
with diagonal \(1+z\), superdiagonal \(z\) and subdiagonal \(1\).
The number of Smith factors of \(H_{n,b}\) having \(p\)-adic
valuation less than \(e\) is
\[
              [z^{b-1}]\bigl(T_u(z)^{n-1}\bigr)_{1,1}.
\]
Indeed the height cutoff from Corollary~\ref{cor:grove-block} is
\(2u-2\). Delete the first and last steps, then read the remaining
steps in pairs between heights \(1,3,\ldots,2u-3\). Marking the
second step of each pair by \(z\) if it rises gives precisely
the entries of \(T_u\); the deleted first step supplies the one
remaining marked rise. In particular \(T_2=(1+z)\), giving the
binary rank in the corollary.

For \(n=5,b=3\), the integral diagonal has six entries \(1\),
thirteen entries \(2\), and one entry \(6\). A direct Smith reduction
of the twenty-by-twenty matrix of joins gives the same list. The
corresponding multiplicities for \(n=8,b=4\) are
\(1^{35},2^{322},6^{126},12^7\).
